\documentclass[10pt,letterpaper]{article}
\usepackage{cornell-notes}

\title{Clause 30.04.07.05: Class Template Moneypunct Byname}
\author{}
\date{}
\setDocTitle{Clause 30.04.07.05: Class Template Moneypunct Byname}
\setDocSubtitle{C++ 2024}
\setDocOwner{C++ 2024 Cornell Notes}
\setDocDate{\today}
\setCornellCollection{C++ 2024 Cornell Notes}
\setCornellUnitType{Clause}
\setCornellUnitNumber{30.04.07.05}
\setCornellUnitTitle{Class Template Moneypunct Byname}
\hypersetup{pdftitle={Clause 30.04.07.05: Class Template Moneypunct Byname},pdfsubject={Cornell notes},pdfauthor={}}

\begin{document}
\maketitle
\begin{CornellOverviewBox}{Overview}
\textbf{Classification:} Standard library - Normative\\[2pt]
\textbf{Learning objective:} Understand the rules, terminology, and semantic consequences associated with Class template moneypunct\_byname.
\end{CornellOverviewBox}\begin{CornellSummaryBox}{Summary}
{\sffamily\bfseries\color{Accent} Summary}\par\vspace{3pt}Section 30.4.7.5 focuses on Class template moneypunct\_byname. Use the listed grammar productions together with the semantic constraints and disambiguation rules referenced by the section. When applying it, preserve the distinction between mandatory requirements, recommendations, permissions, and behavior deliberately left undefined, unspecified, or implementation-defined.
\end{CornellSummaryBox}

\end{document}
